% Copied from chapters/06-queues.tex (fig:queue-fifo).
% Keep the chapter drawing, type sizes, colours, and node styles unchanged.
\begin{utpcompactlabfigure}[The two ends of a queue]
\begin{labgraphic}
\begin{tikzpicture}[
  queueout/.style={draw=UTPTeal,fill=UTPPaleTeal,rounded corners=1mm,
    align=center,font=\scriptsize\color{UTPNavy},inner xsep=2mm,inner ysep=1.2mm},
  queuein/.style={draw=UTPOrange,fill=UTPPaleOrange,rounded corners=1mm,
    align=center,font=\scriptsize\color{UTPNavy},inner xsep=2mm,inner ysep=1.2mm}
]
\foreach \v [count=\i from 0] in {A,B,C,D}{
  \node[utpcell,minimum width=1.5cm,minimum height=0.82cm] (c\i) at (1.5*\i,0) {\v};
}
\node[utptag,anchor=north] at (0,-0.6) {FRONT};
\node[utptag,anchor=north] at (4.5,-0.6) {REAR};
\node[queueout,anchor=east] (deq) at (-1.2,0)
  {\textbf{dequeue}\\[-0.15em]{\tiny oldest leaves}};
\draw[utparrow] (c0.west) -- (deq.east);
\node[queuein,anchor=west] (enq) at (5.7,0)
  {\textbf{enqueue}\\[-0.15em]{\tiny newest joins}};
\draw[utparrow] (enq.west) -- (c3.east);
\end{tikzpicture}
\end{labgraphic}
\end{utpcompactlabfigure}
